%
% Documento: Trabalhos relacionados
%

\chapter{Trabalhos relacionados}
\label{cap:trabalhos-relacionados}

Este capítulo compara estudos que empregam RFID UHF em inventário, rastreamento de equipamentos e pontos de passagem. O objetivo não é reproduzir arquiteturas completas de software, mas identificar quais aspectos devem ser considerados na análise experimental e posteriormente avaliados no protótipo físico.

\section{Inventário e rastreamento de ativos}

\citeonline{chen2022inventory} desenvolveram um sistema de inventário baseado em RFID que integra identificação dos itens e gerenciamento dos registros. O trabalho demonstra que o benefício operacional depende da ligação entre a leitura física e a informação de inventário. Para o presente TCC, essa arquitetura representa o sistema no qual o protótipo poderá ser inserido, mas a camada de aplicação permanece fora do escopo.

\citeonline{ibrahim2018lab} apresentou uma solução de baixo custo para identificação e rastreamento de equipamentos de laboratório com IoT e RFID UHF. Por tratar de equipamentos de laboratório com hardware acessível, esse trabalho é o mais próximo, em finalidade institucional, do cenário de inventário patrimonial universitário abordado aqui. A diferença está no recorte: o presente estudo caracteriza leitor, antena, tags e aquisição serial antes de integrar serviços de rede.

\citeonline{iorga2026} avaliaram o emprego de RFID UHF no inventário automático de componentes por meio de uma bancada com coordenadas definidas. O protocolo mapeou o campo de leitura e examinou histerese, distância, orientação e configurações com uma ou mais tags. Esse controle geométrico é adequado à caracterização espacial do campo. A presente campanha adotou um recorte complementar: observar o desempenho operacional em uma sala institucional comum, sem coordenadas espaciais completas. Essa escolha aproxima o ensaio da aplicação pretendida, embora não permita reproduzir exatamente a geometria nem mapear o canal ponto a ponto.

\section{Arquiteturas centralizadas e segurança}

\citeonline{zohaib2016centralized} propuseram uma arquitetura centralizada para inventário RFID. Já \citeonline{henaien2024smartinventory} discutiram uma arquitetura RFID/IoT com integração de segurança. Esses estudos incluem banco de dados, comunicação e aplicação, enquanto este trabalho isola a etapa de aquisição. A separação é necessária porque uma interface de software bem construída não corrige perdas sistemáticas de leitura.

\section{Cobertura e pontos de passagem}

\citeonline{sabesan2014diversity} mostraram que diversidade de antenas, fase e frequência pode reduzir nulos de multipercurso em áreas amplas. Esse resultado é relevante para salas com superfícies refletoras, mas a solução utiliza múltiplas antenas e supera a capacidade direta do YPD-R200 de uma porta.

Para detectar sentido de deslocamento, \citeonline{wang2022rfaccess} utilizaram duas antenas e um modelo temporal. O trabalho confirma que a identificação de uma tag em uma porta e a inferência de entrada ou saída são problemas distintos. No presente projeto, pontos fixos são discutidos como extensão e não como funcionalidade já entregue pelo protótipo de uma antena.

\section{Síntese comparativa}

A \autoref{tab:trabalhos-relacionados} sintetiza a relação dos estudos com esta pesquisa.

\begin{table}[htbp]
    \centering
    \caption{Comparação entre trabalhos relacionados e o protótipo proposto.}
    \label{tab:trabalhos-relacionados}
    \footnotesize
    \setlength{\tabcolsep}{3pt}
    \renewcommand{\arraystretch}{0.96}
    \begin{tabularx}{\textwidth}{>{\raggedright\arraybackslash}p{2.6cm}>{\raggedright\arraybackslash}p{3.1cm}>{\raggedright\arraybackslash}p{3.0cm}>{\raggedright\arraybackslash}X}
        \toprule
        \textbf{Trabalho} & \textbf{Evidência principal} & \textbf{Variáveis ou métricas} & \textbf{Comparabilidade com este TCC} \\
        \midrule
        \citeonline{chen2022inventory} & Demonstração de sistema integrado de inventário & Identificação e gestão dos registros & Compara-se quanto à finalidade, mas não isola o enlace entre leitor e tag como feito nesta campanha. \\
        \citeonline{ibrahim2018lab} & Protótipo de baixo custo em laboratório & Identificação e rastreamento de equipamentos & É próximo no contexto institucional; a integração IoT excede o escopo deste protótipo. \\
        \citeonline{iorga2026} & Bancada controlada e cinco sequências de teste & Campo de leitura, histerese, distância e orientação & É a referência metodológica mais próxima, embora utilize objetos, trajetória e hardware distintos. \\
        \citeonline{sabesan2014diversity} & Ensaio de cobertura com diversidade & Ganho de cobertura por antena, fase e frequência & Oferece referência metodológica para multipercurso, mas usa arquitetura multiaxial incompatível com uma única porta. \\
        \citeonline{wang2022rfaccess} & Experimento de controle de passagem & Sequência temporal em duas antenas & Fundamenta requisitos de direção; não é quantitativamente comparável ao ensaio de presença em uma zona. \\
        \bottomrule
    \end{tabularx}
    \fonte{Elaborado pelo autor a partir das referências citadas (2026).}
\end{table}

Os estudos variam em ambiente, hardware e nível de integração; por isso, seus resultados não devem ser comparados diretamente por percentuais com a campanha deste TCC. A síntese mostra, porém, critérios metodológicos comuns: registrar a geometria, separar cobertura de simples contagem de quadros e distinguir detecção de presença de inferência de direção.

Não foi identificado, entre as referências acadêmicas selecionadas, um estudo revisado por pares que caracterize exatamente a variante YPD-R200, a antena APCA8090 e a etiqueta AD-239 M730 utilizadas aqui. Para o módulo, as fontes específicas disponíveis são o datasheet, o protocolo serial do fabricante e a implementação pública de \mbox{\citeonline{playfultechnology2023r200}}. Como a designação comercial ``R200'' aparece em produtos diferentes, resultados obtidos com outros leitores de mesmo nome não foram tratados como equivalentes. Essa lacuna reforça a contribuição de documentar a configuração do protótipo, o firmware desenvolvido e os resultados experimentais obtidos com o conjunto efetivamente ensaiado.

O diferencial deste trabalho não é propor uma nova interface aérea nem um sistema patrimonial completo. A contribuição está em adaptar o firmware de um leitor YPD-R200 de baixo custo para Arduino e analisar experimentalmente condições relevantes de antena, etiqueta, material e aquisição de EPCs. Esse recorte estrutura hipóteses, métricas e critérios de comparação para uma evolução futura do inventário ou de portais fixos.
